\pdfinfoomitdate=1
\pdftrailerid{}
\pdfsuppressptexinfo=15
\documentclass[11pt]{article}
\usepackage[T1]{fontenc}
\usepackage{lmodern}
\usepackage{amsmath,amssymb,amsthm}
\usepackage{booktabs}
\usepackage{geometry}
\usepackage[hidelinks]{hyperref}
\usepackage{microtype}
\geometry{margin=1in}

\newtheorem{theorem}{Theorem}[section]
\newtheorem{lemma}[theorem]{Lemma}
\newtheorem{proposition}[theorem]{Proposition}
\newtheorem{corollary}[theorem]{Corollary}
\theoremstyle{definition}
\newtheorem{definition}[theorem]{Definition}
\newtheorem{computation}[theorem]{Exact computation}
\theoremstyle{remark}
\newtheorem{remark}[theorem]{Remark}

\newcommand{\QQ}{\mathbf Q}
\newcommand{\PP}{\mathbf P}
\newcommand{\Spec}{\operatorname{Spec}}
\newcommand{\Proj}{\operatorname{Proj}}
\newcommand{\im}{\operatorname{im}}
\newcommand{\rank}{\operatorname{rank}}
\newcommand{\Der}{\operatorname{Der}}
\newcommand{\Ext}{\operatorname{Ext}}
\newcommand{\Gr}{\operatorname{Gr}}
\newcommand{\ad}{\operatorname{ad}}

\title{A smoothing of the Gr\"unbaum Stanley--Reisner threefold\\
\large Standalone proof and exact computational certificate}
\author{}
\date{}

\begin{document}
\maketitle

\begin{abstract}
We prove that the Stanley--Reisner threefold of the Gr\"unbaum--Sreedharan
eight-vertex triangulated $3$-sphere is the special fibre of a flat projective
family over $\QQ[[s]]$ with smooth geometric generic fibre.  The proof has
three logically separate parts.  First, an intrinsic Andr\'e--Quillen bracket
calculation and the curved Bianchi identity construct an all-orders formal
deformation.  Second, 280 exact singular-incidence calculations show that its
rational two-jet already forces generic smoothness, independently of all
higher coefficients.  Third, Grothendieck existence algebraizes the formal
closed subscheme, and a torsion argument proves flatness.  All
$L_\infty$-algebras below use the ordinary unshifted cohomological convention:
$\ell_n$ has degree $2-n$.  Every computer-assisted assertion is identified
by its exact input, algorithm, output, and mathematical consequence.
\end{abstract}

\tableofcontents

\section{The special fibre and the computational contract}

Put $k=\QQ$, $S=k[x_1,\ldots,x_8]$, and let $\Delta$ have the following
twenty tetrahedral facets:
\[
\begin{gathered}
1234,1237,1267,1347,1567,2345,2367,3467,3456,4567,\\
2358,2368,3568,1268,1568,1248,2458,1478,1578,4578.
\end{gathered}
\]
This is the Gr\"unbaum--Sreedharan triangulation in the labelling fixed in
\cite{GS}.  Its minimal nonfaces are the sixteen triples occurring in
\[
\begin{split}
I=(\,&x_6x_7x_8,x_4x_6x_8,x_3x_7x_8,x_3x_5x_7,
x_3x_4x_8,x_2x_7x_8,\\
&x_2x_5x_7,x_2x_5x_6,x_2x_4x_7,x_2x_4x_6,
x_1x_4x_6,x_1x_4x_5,\\
&x_1x_3x_8,x_1x_3x_6,x_1x_3x_5,x_1x_2x_5\,)\subset S.
\end{split}                                                    \tag{1.1}
\]
Denote the ordered cubic generators displayed in (1.1) by
$m_0,\ldots,m_{15}$.  Write $A=S/I$ and $X_0=\Proj A\subset\PP^7_k$.

\begin{computation}[special-fibre combinatorics]\label{comp:comb}
The input of \path{code/verify_combinatorics.py} is exactly the twenty
lines of \path{data/grunbaum_facets.txt}.  The script enumerates all faces
and all inclusion-minimal nonfaces.  For every face $\sigma$, including the
empty face, it constructs the simplicial boundary matrices of
$\operatorname{lk}_\Delta(\sigma)$ and row-reduces them over $\QQ$ using
\texttt{Fraction} arithmetic.  Thus no floating-point calculation occurs.

The output is: 8 vertices, 20 facets, 16 minimal nonfaces (precisely the
triples in (1.1)), 97 faces in total, pure dimension 3, every ridge in exactly
two facets, connected facet-dual graph, and
\[
 \widetilde b(\Delta;\QQ)=(0,0,0,1),\qquad
 \widetilde H_i(\operatorname{lk}_\Delta\sigma;\QQ)=0
 \quad(i<\dim\operatorname{lk}_\Delta\sigma)
\]
for every $\sigma$.  By Reisner's criterion \cite{Reisner}, $A$ is
Cohen--Macaulay of Krull dimension $4$.  Consequently $X_0$ is a pure
Cohen--Macaulay projective threefold.  The certificate is
\path{verification/combinatorics_QQ.txt}.
\end{computation}

We use intrinsic Andr\'e--Quillen cohomology.  Let $P\to A$ be a semi-free
commutative Tate differential graded $k$-algebra resolution, homologically
graded with differential of degree $-1$.  A derivation of cohomological degree
$p$ lowers homological degree by $p$.  Set
\[
 C^p=\Der_k^p(P,A)_0,\qquad T^p_0(A)=H^p(C^\bullet),             \tag{1.2}
\]
where the subscript denotes internal degree zero.  Formula (1.2), rather than
a naive higher Ext group of $I/I^2$, is the definition used below.

\begin{remark}[an intentionally absent claim]\label{rem:no-t3dim}
We do \emph{not} claim that $\dim T^3_0=24$.  A computation of
$\Ext_A^2(I/I^2,A)_0$ gives 24, but for a non-lci quotient that group is not
automatically intrinsic $T^3$.  The proof needs only the twelve independent
intrinsic $T^3$ classes established in Proposition~\ref{prop:bracket}.
\end{remark}

\section{Construction of the all-orders formal deformation}

\subsection{The intrinsic tangent Lie algebra and the unshifted convention}

The dg Lie algebra $\mathfrak g=\Der_k(P,P)_0$, with differential
$[d,-]$ and graded commutator, controls graded deformations of $A$; this is the
tangent Lie algebra construction for a cofibrant resolution
\cite{Hinich97,Hinich04}.  Since $\Omega_{P/k}$ is cofibrant as a dg
$P$-module, applying $\operatorname{Hom}_P(\Omega_{P/k},-)$ to the
quasi-isomorphism $P\to A$ gives a quasi-isomorphism
\[
 \Der_k(P,P)_0\longrightarrow\Der_k(P,A)_0.                    \tag{2.1}
\]
Thus the bracket induced on the cohomology of (1.2) is the intrinsic
Andr\'e--Quillen tangent-Lie bracket.  It is not Yoneda multiplication on a
conormal Ext group.

Transfer $\mathfrak g$ to its cohomology $H$.  We use the ordinary
cohomological $L_\infty$ convention throughout:
\[
 \ell_n:\bigwedge^n H\longrightarrow H,\qquad |\ell_n|=2-n.
\]
The model is minimal, so $\ell_1=0$.  For $t\in H^1\widehat\otimes\mathfrak m$
define
\[
 \kappa(t)=\sum_{n\ge2}\frac1{n!}\ell_n(t,\ldots,t)\in
 H^2\widehat\otimes\mathfrak m^2,                              \tag{2.2}
\]
and, for $w\in H^2$,
\[
 d_t(w)=\sum_{n\ge0}\frac1{n!}
 \ell_{n+1}(\underbrace{t,\ldots,t}_{n},w)\in H^3.             \tag{2.3}
\]
The $L_\infty$ identities, applied to the curvature (2.2), give the curved
Bianchi identity
\[
 d_t\kappa(t)=0.                                                \tag{2.4}
\]
Notice that (2.2)--(2.4) are unshifted formulas: Maurer--Cartan elements have
degree $1$, obstructions degree $2$, and Bianchi values degree $3$.

The completed Maurer--Cartan base is a hull for the graded deformation
functor.  Make the change of a minimal hull presentation explicit.  If the
27 equations in another set of 53 tangent coordinates are
\[
 G(t)=U(t)\,\kappa(\phi(t)),
\]
where $\phi$ has invertible linear part $\phi_1$ and $U(0)$ is invertible,
then
\[
 B_G(t)=d_{\phi(t)}\circ U(t)^{-1}
 \quad\hbox{satisfies}\quad B_G(t)G(t)=0.
\]
Its term linear in $t$ is
$\ad_{\phi_1(t)}\circ U(0)^{-1}$.  We henceforth use the coordinates
produced by the pinned \texttt{VersalDeformations} computation.  The
constant relation-basis transport in (2.9) identifies its 27 obstruction
coordinates with the Tate model, while the direct primary-commutator
comparison in Proposition~\ref{prop:bracket} checks that the linear operator
in the preceding display, in exactly these coordinates, is the computed
intrinsic bracket
(and that the differential of the quadratic equations is $-1$ times the
primary commutator).  Thus the Bianchi induction below is invariant under
the presentation change, rather than being an identity assumed of package
coordinates.  From now on we again write $\kappa$ and $d_t$ for this
transported tuple $G$ and operator $B_G$; (2.4) denotes the transported
identity as well.

\subsection{The rational direction and its initial two-jet}

In the computed basis of $T^1_0(A)$, take the explicit vector
\[
\begin{split}
y=(\,&1,3,1,2,6,1,1,1,1,1,21,28,1,5,10,55,66,15,1,18,2,\\
&35,42,1,4,10,1,25,12,30,5,33,44,1,6,3,14,7,1,9,1,6,\\
&12,20,1,4,22,1,15,1,1,8,11\,).
\end{split}                                                     \tag{2.5}
\]
No property of how (2.5) was found is used.

\begin{computation}[the degree-three miniversal calculation]
\label{comp:formaldata}
The exact inputs of \path{code/verify_formal_data.m2} are (1.1) and
(2.5).  Over $\QQ$, \texttt{CT\^1(0,F0)} and \texttt{CT\^2(0,F0)} compute
the degree-zero tangent and second cotangent-cohomology bases.  The command
\texttt{versalDeformation} is run through order three with
\texttt{SmartLift=>false} and \texttt{DefParam=>t}.  Its algorithm solves the
deformation equation and records its obstruction coordinates by exact
polynomial and module arithmetic.

The asserted output is
\[
 \dim T^1_0=53,\quad \dim T^2_0=27,\quad q(y)=0,\quad
 \rank Dq_y=15,\quad \kappa_3(y)=0,                            \tag{2.6}
\]
where $q=\kappa_2$ in these coordinates and $\kappa_3$ is the cubic base
term.  The 27 quadratic equations are also checked to be a minimal tuple.
The same run specializes the miniversal family along $t=sy$ and writes its
sixteen first- and second-order cubic corrections to
\path{data/universal_2jet_QQ.txt}; its SHA-256 digest is
\[
\texttt{40e64e61674b6a4e61f1ea6822dc79327bf4ba397285f84ed8738c4c18cd1795}.
                                                                    \tag{2.7}
\]
Thus $t_2(s)=sy$ has zero curvature through $s^3$ and gives the displayed
family two-jet.  The certificate is
\path{verification/formal_data_QQ.txt}.
\end{computation}

\subsection{The intrinsic bracket calculation}

We now explain both what the computer calculates and why it is intrinsic.
The materialized part of the semi-free resolution has:
\[
\begin{array}{c|c|c}
\text{generators}&\text{homological degree}&\text{internal degree}\\ \hline
e_1,\ldots,e_{16}&1&3\\
f_1,\ldots,f_{30}&2&4\\
g_1,\ldots,g_{136}&3&5\text{ (16 times), }6\text{ (120 times).}
\end{array}                                                       \tag{2.8}
\]
Here $d(e_i)$ is the $i$th generator of (1.1); the $f_j$ kill a basis of
first syzygies, and the $g_a$ kill all cycles at the next stage.

Positive-degree derivations vanish on the polynomial subring $S$ for degree
reasons.  The degree-two cochain differential is therefore the linear
$f$-part of the 136 elements $d(g_a)$.  The source formula used by
\texttt{CT\^2},
\[
 \frac{\operatorname{Hom}_A((\operatorname{syz}F_0)/
        (\text{Koszul syzygies}),A)}
      {\im(\,{}^{\mathsf T}\!\operatorname{syz}F_0\,)},         \tag{2.9}
\]
is exactly $\ker\delta_2/\im\delta_1=T^2_0(A)$.  The package syzygy basis
and the $f$-basis in (2.8) are related by a $30\times30$ matrix of constants.
The script solves for it, checks its determinant is a unit, transports all 27
classes by its transpose, and checks their closure on all 136 $g_a$.

Let $\theta$ be the degree-one derivation representing $y$, and let
$\eta_1,\ldots,\eta_{27}$ represent the transported $T^2_0$ basis.  The
script first solves the closure equations for their values on every generator
in (2.8), then computes
\[
 [\theta,\eta_j](g_a)=
 \theta(\eta_j(g_a))-\eta_j(\theta(g_a)).                       \tag{2.10}
\]
It finally reduces the 27 columns (2.10) modulo the complete image of
$\delta_2:C^2_0\to C^3_0$.

The following extension argument is essential; without it a calculation on a
truncated resolution would not prove a statement in $T^3$.

\begin{lemma}[Tate extension]\label{lem:tate}
A degree-$p$ derivation, $p=1$ or $2$, which is closed on the generators of
(2.8), extends without changing those values to a closed derivation of the
full Tate resolution $P$.
\end{lemma}

\begin{proof}
Proceed through the remaining Tate generators in increasing homological
degree.  If $z$ is newly adjoined in degree $n\ge4$, closure prescribes
\[
 d(\vartheta z)=(-1)^p\vartheta(dz).                            \tag{2.11}
\]
Indeed, closedness $[d,\vartheta]=0$ is exactly (2.11).  Because $d^2z=0$ and
$\vartheta$ is already closed on earlier generators, the right side of
(2.11) is a cycle of homological degree $n-p-1$.  This degree is positive.
The Tate resolution is acyclic in positive homological degrees, so the cycle
is a boundary.  A primitive has degree $n-p<n$, hence uses only generators
already constructed.  The differential preserves internal degree, so the
primitive can be chosen in the required internal homogeneous component.
This defines $\vartheta(z)$ and completes the induction.  The low-degree
cases, where positivity alone would not suffice, are precisely the equations
solved and asserted by the script on (2.8).
\end{proof}

By Lemma~\ref{lem:tate}, $\theta$ and each $\eta_j$ extend globally.  Their
commutators are therefore globally closed.  Consequently the columns (2.10)
lie in $\ker\delta_3$, and the natural map
\[
 T^3_0(A)=\ker\delta_3/\im\delta_2
 \longrightarrow C^3_0/\im\delta_2                             \tag{2.12}
\]
is injective on their span.  Rank computed in the right side of (2.12) is
therefore their actual rank in intrinsic Andr\'e--Quillen $T^3_0$.

\begin{proposition}[fixed-direction exactness]\label{prop:bracket}
For the direction (2.5), the intrinsic bracket satisfies
\[
 \rank\bigl(\ad_y:T^2_0(A)\to T^3_0(A)\bigr)=12,\qquad
 \ker\ad_y=\im Dq_y.                                           \tag{2.13}
\]
\end{proposition}

\begin{proof}
The exact inputs of \path{code/verify_aq_bracket.m2} are again (1.1) and
(2.5).  The algorithm constructs (2.8), performs the basis transport (2.9),
solves all closure equations, forms (2.10), and uses exact Gr\"obner/module
arithmetic over $\QQ$ to compute its rank in (2.12).  Its output is rank 12.

As an independent coordinate check inside the same run, the primary
commutator of $\theta_y$ with each of the 53 degree-one basis derivations is
computed in $T^2_0$.  Its matrix is exactly $-Dq_y$ in the obstruction
coordinates returned by \texttt{VersalDeformations}; the minus sign is the
differential convention and does not change its image.  The script also
checks directly that
\[
 \ad_yDq_y=0,\qquad \rank Dq_y=15.                              \tag{2.14}
\]
Thus $\im Dq_y\subseteq\ker\ad_y$.  The domain of $\ad_y$ has dimension 27,
so its kernel has dimension $27-12=15$, equal to the dimension of
$\im Dq_y$.  This proves (2.13).  The complete output is
\path{verification/aq_bracket_QQ.txt}.
\end{proof}

\subsection{Bianchi induction to every order}

\begin{proposition}[all-orders arc]\label{prop:arc}
There is a series
\[
 t(s)=sy+v_2s^2+v_3s^3+\cdots\in T^1_0(A)[[s]],\qquad v_2=0,    \tag{2.15}
\]
over $\QQ$ whose curvature vanishes identically and whose induced embedded
family has the two-jet (2.7).
\end{proposition}

\begin{proof}
By (2.6), $t_2(s)=sy$ has curvature $O(s^4)$.  Suppose inductively that a
partial series has been chosen with
\[
 \kappa(t(s))=r_Ns^N+O(s^{N+1}),\qquad N\ge4.                  \tag{2.16}
\]
In the coefficient of $s^{N+1}$ in (2.4), all terms involving a curvature
coefficient below $N$ vanish.  Since $t(s)=sy+O(s^2)$ and $\ell_1=0$, the
only remaining term is
\[
 \ell_2(y,r_N)=\ad_y(r_N)=0.                                  \tag{2.17}
\]
Proposition~\ref{prop:bracket} gives $r_N\in\im Dq_y$.  Choose
$v_{N-1}\in T^1_0(A)$ with
\[
 Dq_y(v_{N-1})=-r_N.                                          \tag{2.18}
\]
Adding $v_{N-1}s^{N-1}$ to $t(s)$ changes the coefficient of $s^N$ in
the quadratic term by $Dq_y(v_{N-1})$.  In an $m$-ary term with $m\ge3$,
one occurrence of the new coefficient has order at least
$(N-1)+(m-1)\ge N+1$; two occurrences have still higher order.  Hence
(2.18) kills precisely the residual in (2.16) without altering lower orders.
Induction constructs (2.15), and completeness gives $\kappa(t(s))=0$.
Every step solves a rational linear system, so all coefficients lie in
$\QQ$.  Pullback of the formal miniversal family produces compatible flat
graded deformations over $\QQ[s]/(s^n)$ for every $n$.
\end{proof}

\begin{lemma}[the formal ideal has sixteen cubic generators]
\label{lem:cubics}
Let $R=\QQ[[s]]$.  The formal embedded family of
Proposition~\ref{prop:arc} is generated by homogeneous cubics
\[
 F_i=m_i+sg_i+s^2h_i+s^3R_i(s,x)\in
 R[x_1,\ldots,x_8]\qquad(0\le i<16).                         \tag{2.19}
\]
Their reductions modulo $s^3$ are the file (2.7).
\end{lemma}

\begin{proof}
Write $\widehat{\mathcal A}=\varprojlim_n\mathcal A_n$ for the compatible
flat graded quotients furnished by the formal miniversal family, and let
$\widehat{\mathcal I}$ be the degreewise kernel of
$R[x_1,\ldots,x_8]\to\widehat{\mathcal A}$.  For each degree $d$, the
module $(\mathcal A_n)_d$ is finite flat, hence free, over the Artinian local
ring $R/(s^n)$, with rank independent of $n$.  Its inverse limit is therefore
a finite free $R$-module.  It follows from
\[
 0\longrightarrow\widehat{\mathcal I}_d
 \longrightarrow R[x_1,\ldots,x_8]_d
 \longrightarrow\widehat{\mathcal A}_d\longrightarrow0
\]
that $\widehat{\mathcal I}_d$ is finite free as well.

The special degree-three kernel has basis $m_0,\ldots,m_{15}$ from (1.1).
The sixteen computed equations modulo $s^3$ lift, by surjectivity from the
finite free module $\widehat{\mathcal I}_3$, to cubics $F_i$ with exactly the
two-jet (2.7).  Their reductions modulo $s$ form the displayed special basis,
so Nakayama's lemma makes them an $R$-basis of
$\widehat{\mathcal I}_3$.  For every
$d\ge3$, the degree-$d$ part of
$\widehat{\mathcal I}/(F_0,\ldots,F_{15})$ is a finite $R$-module whose
reduction modulo $s$ is zero, because the special ideal (1.1) is generated
by its cubics.  Nakayama again makes this quotient zero.  For $d<3$, the same
argument applied directly to $\widehat{\mathcal I}_d$ gives zero.  Hence the sixteen
cubics generate the full formal ideal.  A degree-three polynomial has only
finitely many $x$-monomials, so its $s$-adic coefficient series make it an
honest element of $R[x_1,\ldots,x_8]$, proving the claim.
\end{proof}

\section{The geometric generic fibre is smooth}

Let
\[
 F_i^{(2)}=m_i+sg_i+s^2h_i\quad(0\le i<16)                    \tag{3.1}
\]
be the cubics in \path{data/universal_2jet_QQ.txt}.  By
Lemma~\ref{lem:cubics}, the family from Proposition~\ref{prop:arc} has equations
\[
 F_i=F_i^{(2)}+s^3R_i(s,x),                                   \tag{3.2}
\]
with $R_i\in\QQ[[s]][x_1,\ldots,x_8]$ homogeneous of degree three in $x$.

On the projective chart $x_v=1$, write $z_1,\ldots,z_7$ for the remaining
coordinates and let $J$ be the $7\times16$ affine Jacobian of the equations.
The flat family is Cohen--Macaulay of relative dimension three near the
special fibre: this follows from Computation~\ref{comp:comb} and the standard
flat Cohen--Macaulay fibre criterion (compare \cite[Tag 0E7T]{Stacks}).
Thus a point of a fibre is singular exactly when $\rank J\le3$.  Equivalently,
$\ker J^{\mathsf T}\subset k^7$ contains a four-plane.

The Grassmannian $\Gr(4,7)$ has $\binom74=35$ standard affine charts.  On the
chart indexed by a four-subset $P\subset\{1,\ldots,7\}$, let $K_P$ be the
$7\times4$ matrix whose rows indexed by $P$ form the identity and whose other
entries are twelve variables $a_1,\ldots,a_{12}$.  For a pair $(v,P)$ the
singular-incidence ideal is
\[
 \mathcal J_{v,P}=(F_0,\ldots,F_{15},\ J^{\mathsf T}K_P).       \tag{3.3}
\]
It has $16+16\cdot4=80$ displayed generators.

\begin{computation}[all 280 incidence memberships]\label{comp:incidence}
The exact input is (3.1), verified by digest (2.7).  For each of the
$8\cdot35=280$ pairs, \path{code/verify_incidence_chart.m2}
dehomogenizes (3.1), constructs $J$ and $K_P$, and truncates (3.3) modulo
$s^3$.  Put
\[
 B=\QQ[z_1,\ldots,z_7,a_1,\ldots,a_{12}].
\]
For every generator $q=q_0+sq_1+s^2q_2$, multiplication by all elements of
$B[s]/(s^3)$ is represented exactly by the three columns
\[
 (q_0,q_1,q_2)^{\mathsf T},\quad(0,q_0,q_1)^{\mathsf T},\quad
 (0,0,q_0)^{\mathsf T}.                                      \tag{3.4}
\]
The 80 generators therefore give a $3\times240$ module matrix.  Membership of
$(0,0,1)^{\mathsf T}$ in its column module is exactly the assertion
\[
 s^2\in(\mathcal J_{v,P},s^3).                                \tag{3.5}
\]

The driver \path{code/verify_all_incidence.py} enumerates every pair and
uses Macaulay2 exact rational module Gr\"obner reduction.  A zero remainder
against a partial Gr\"obner basis is already a proof of membership: each
basis column is obtained by exact module operations from the 240 input
columns, and the zero reduction expresses the target in those columns.  A
nonzero remainder is treated only as inconclusive.  Degree six proves 277
pairs; degree eight proves the remaining pairs
\[
 (v,P)=(2,7),(6,12),(8,16).                                   \tag{3.6}
\]
The driver checks that its sorted result contains all 280 pairs.  On pair
$(1,1)$ it additionally retains the change matrix and verifies the resulting
identity in the original columns (eight coefficients are nonzero).  The exact
output is \path{verification/incidence_all_QQ.txt}.  The mathematical
consequence is (3.5) on every incidence chart.
\end{computation}

We spell out why 280 memberships exhaust every possible singular generic
point and why the truncation is uniform in (3.2).

\begin{proposition}[two-jet smoothness]\label{prop:smooth}
Every flat projective continuation (3.2) has smooth geometric generic fibre.
\end{proposition}

\begin{proof}
Suppose a geometric generic singular point exists.  The singular-incidence
scheme is of finite type over $K=\QQ((s))$, so it has a point after some finite
extension $K'/K$.  Choose a DVR $R'$ in the integral closure of $\QQ[[s]]$ in
$K'$ with fraction field $K'$.  Properness of $\PP^7$ extends the projective
point to $\Spec R'$.  Scale its homogeneous coordinates so that they are all
integral and at least one, say $x_v$, is a unit.  This puts the section in one
of the eight affine projective charts.

At the generic point, singularity means $\dim\ker J^{\mathsf T}\ge4$; choose
a four-dimensional subspace.  It is a $K'$-point of the proper Grassmannian
$\Gr(4,7)$ and hence extends over $R'$.  Its Pl\"ucker coordinates may be
scaled to be integral with at least one unit.  The unit Pl\"ucker coordinate
selects one of the 35 standard charts, and the twelve entries of $K_P$ on
that chart are integral.  Thus some one of the 280 pairs contains integral
coordinates for both the point and its kernel four-plane.  The equations
$J^{\mathsf T}K_P=0$, true over $K'$, remain true over $R'$ because their
entries are integral and $R'$ injects into $K'$.

Fix this chart.  Membership (3.5), lifted from the truncated polynomial ring,
is an identity
\[
 s^2=\sum_{j=1}^{80}c_jq_j+s^3c_0                            \tag{3.7}
\]
with $c_j,c_0$ polynomial in $s,z,a$ and $q_j$ the truncated incidence
generators.  For the actual continuation, each generator is
$q'_j=q_j+s^3r_j$: differentiation with respect to an affine coordinate does
not lower $s$-adic order.  Substitution in (3.7) gives
\[
 s^2=\sum_j c_jq'_j+s^3\left(c_0-\sum_jc_jr_j\right).          \tag{3.8}
\]
All quantities in parentheses evaluate to elements of $R'$ because every
projective and Grassmann coordinate is integral and the coefficients of
(3.2) lie in $\QQ[[s]]$.  Along the incidence section all $q'_j$ vanish, so
(3.8) becomes
\[
 s^2=s^3c\qquad(c\in R').                                     \tag{3.9}
\]
If $v$ is the valuation of $R'$, then $v(s)>0$, whereas (3.9) implies
$2v(s)=3v(s)+v(c)\ge3v(s)$, a contradiction.  Hence no geometric generic
singular point exists.
\end{proof}

\section{Algebraization and flatness}

Set $R=\QQ[[s]]$ and $R_n=R/(s^n)$.  Proposition~\ref{prop:arc} gives
cartesian closed subschemes
\[
 X_n\subset\PP^7_{R_n},\qquad X_{n+1}\times_{R_{n+1}}R_n=X_n, \tag{4.1}
\]
each flat over $R_n$, with $X_1=X_0$.  We now justify both effectivity and
flatness of the algebraized family.

\begin{proposition}[effectivity]\label{prop:effectivity}
The system (4.1) is the system of reductions of a closed subscheme
$X\subset\PP^7_R$ that is proper over $R$.
\end{proposition}

\begin{proof}
The ring $R$ is Noetherian and complete for $(s)$, and $\PP^7_R$ is separated
and of finite type (indeed proper) over $R$.  The special member $X_1$ is
projective over $\QQ$.  Thus every hypothesis of Grothendieck's algebraization
theorem for compatible closed subschemes is satisfied; in the formulation
\cite[Tag 0899]{Stacks}, (4.1) algebraizes to a closed immersion
$X\hookrightarrow\PP^7_R$, and $X$ is proper over $R$.

For completeness, the theorem applies to the compatible coherent quotients
$\mathcal O_{\PP^7_{R_n}}\twoheadrightarrow\mathcal O_{X_n}$.
Grothendieck existence algebraizes the quotient and its map; exactness and
full faithfulness of completion show that the algebraized map is still
surjective, hence its kernel is a coherent ideal sheaf.  This is why existence
of a coherent sheaf alone is not being mistaken for existence of a closed
subscheme.

There is also a concrete description.  Lemma~\ref{lem:cubics} gives sixteen
homogeneous cubics in $R[x_1,\ldots,x_8]$.  Their Proj has reductions (4.1), and
the full-faithfulness part of Grothendieck existence identifies it with the
closed subscheme just algebraized.  Consequently the equations used in the
DVR argument are equations of the algebraized family, not merely unrelated
formal functions.
\end{proof}

\begin{proposition}[flatness]\label{prop:flat}
The morphism $X\to\Spec R$ is flat.
\end{proposition}

\begin{proof}
It is enough to prove $R$-flatness at every stalk $B=\mathcal O_{X,x}$.
If $x$ is on the generic fibre, $s$ is a unit.  Suppose $x$ specializes to
the closed fibre and $sb=0$ in $B$.  For every $n\ge2$, the reduction
$B/s^nB$ is flat over $R_n$ by (4.1).  In $R_n$ the annihilator of $s$ is
$(s^{n-1})$.  Tensoring the corresponding kernel sequence with the flat
$R_n$-module $B/s^nB$ gives
\[
 b\bmod s^nB\in s^{n-1}(B/s^nB).
\]
Thus $b\in s^{n-1}B$ for every $n$.  The local ring $B$ is Noetherian and
$s$ lies in its Jacobson radical, so Krull intersection gives
$\bigcap_ns^nB=0$ and hence $b=0$.  Therefore $B$ is $s$-torsion-free.
Since $R$ is a DVR, an $R$-module is flat exactly when it is torsion-free.
Thus every stalk is $R$-flat.
\end{proof}

\begin{theorem}[smoothing theorem]\label{thm:main}
The Gr\"unbaum Stanley--Reisner threefold $X_0$ is smoothable in
characteristic zero.  More precisely, there is a flat projective morphism
\[
 X\longrightarrow\Spec\QQ[[s]]
\]
whose special fibre is $X_0$ and whose geometric generic fibre is smooth.
\end{theorem}

\begin{proof}
Proposition~\ref{prop:arc} constructs the compatible formal family;
Propositions~\ref{prop:effectivity} and \ref{prop:flat} algebraize it and prove
flatness; Proposition~\ref{prop:smooth} proves smoothness after algebraic
closure of the fraction field.  The special fibre is (1.1) by construction.
\end{proof}

\section{Audit table for computer-assisted lemmas}

\begin{center}
\begin{tabular}{p{0.18\textwidth}p{0.25\textwidth}p{0.25\textwidth}p{0.21\textwidth}}
\toprule
Claim & Exact input & Exact algorithm and output & Consequence \\
\midrule
Special fibre & Twenty facets in \path{data/} & Enumerate faces; rational boundary-matrix ranks; 97 Reisner checks & Pure CM threefold; ideal (1.1) \\
Initial arc data & Ideal (1.1), vector (2.5) & M2 miniversal deformation to degree 3; (2.6), digest (2.7) & Initial zero residual and fixed two-jet \\
Intrinsic bracket & Same ideal and vector & Three Tate stages; extendable derivations; commutators modulo all $\delta_2$ boundaries; ranks 12 and 15 & $\ker\ad_y=\im Dq_y$ \\
Generic smoothness & The regenerated two-jet & 280 exact module memberships over $\QQ$; 277 at degree 6, 3 at degree 8 & DVR contradiction for every higher continuation \\
\bottomrule
\end{tabular}
\end{center}

All four rows are run by \path{./verify.sh}.  The scripts contain assertions
for dimensions, basis changes, closure equations, ranks, chart coverage,
semantic output, and hashes.  The calculations are indispensable parts of
the proof and are not concealed as heuristic evidence.

\begin{thebibliography}{9}

\bibitem{GS}
B.~Gr\"unbaum and V.~P.~Sreedharan,
\emph{An enumeration of simplicial $4$-polytopes with $8$ vertices},
J. Combinatorial Theory \textbf{2} (1967), 437--465,
\href{https://doi.org/10.1016/S0021-9800(67)80055-3}{doi:10.1016/S0021-9800(67)80055-3}.

\bibitem{Hinich97}
V.~Hinich,
\emph{Homological algebra of homotopy algebras},
Comm. Algebra \textbf{25} (1997), 3291--3323,
\href{https://arxiv.org/abs/q-alg/9702015}{arXiv:q-alg/9702015}.

\bibitem{Hinich04}
V.~Hinich,
\emph{Deformations of homotopy algebras},
Comm. Algebra \textbf{32} (2004), 473--494,
\href{https://doi.org/10.1081/AGB-120027907}{doi:10.1081/AGB-120027907}.

\bibitem{Ilten}
N.~O.~Ilten,
\emph{Versal deformations and local Hilbert schemes},
J. Software Algebra Geom. \textbf{4} (2012), 12--16,
\href{https://doi.org/10.2140/jsag.2012.4.12}{doi:10.2140/jsag.2012.4.12}.

\bibitem{Reisner}
G.~A.~Reisner,
\emph{Cohen--Macaulay quotients of polynomial rings},
Adv. Math. \textbf{21} (1976), 30--49,
\href{https://doi.org/10.1016/0001-8708(76)90114-6}{doi:10.1016/0001-8708(76)90114-6}.

\bibitem{Stacks}
The Stacks Project Authors,
\emph{The Stacks Project},
\href{https://stacks.math.columbia.edu/tag/0899}{Tag 0899}
and \href{https://stacks.math.columbia.edu/tag/0E7T}{Tag 0E7T}.

\end{thebibliography}

\end{document}
